% =====================================================================
% When the Benchmark Grades Itself: Two Information Leaks in
% Football Player-Representation Evaluation
%
% PORTING NOTE (ECIR / Springer LNCS):
% This file uses the standard `article` class for a clean, self-contained
% build. To port to Springer LNCS (e.g. for an ECIR submission):
%   1. Replace the preamble block marked "ARTICLE PREAMBLE" with:
%        \documentclass{llncs}
%        \usepackage{graphicx,booktabs,amsmath,amssymb,url}
%   2. Replace \maketitle's \author/\date block with the llncs
%        \author{...}\institute{...} form, and add \maketitle.
%   3. Replace the unnumbered abstract environment with the llncs
%        \begin{abstract}...\end{abstract} placed before \section{Introduction}.
%   4. The body (\section/\subsection, tables, \cite, bibliography) is
%        already LNCS-compatible and needs no further change.
% =====================================================================

% --------------------------- ARTICLE PREAMBLE ------------------------
\documentclass[11pt]{article}
\usepackage[margin=1in]{geometry}
\usepackage{graphicx}
\usepackage{booktabs}
\usepackage{amsmath}
\usepackage{amssymb}
\usepackage{url}
\usepackage[hidelinks]{hyperref}
\usepackage{microtype}
\usepackage{enumitem}

\title{When the Benchmark Grades Itself:\\
Two Information Leaks in Football\\ Player-Representation Evaluation}
\author{Aditya Patni\\ Independent Researcher\\ \texttt{aditya@latimal.com}}
\date{June 2026}
% ---------------------------------------------------------------------

\begin{document}
\maketitle

\begin{abstract}
Learned representations of football (soccer) players are routinely reported to
capture playing style, yet the field validates them in ways that leak the
relevance signal into the model being graded. We identify two such leaks, give
leakage-free protocols, and release an open benchmark and evaluator. First, a
set-transformer over anonymized StatsBomb 360 freeze-frames scores 0.878
within-tournament identity Top-1 under a random split but collapses to 0.188
under a match-disjoint split (disjoint confidence intervals), with a same-match
control at 0.889 isolating match-context rather than enduring style; the
collapse reproduces for static heatmap representations (a parameter-free heatmap
and a Player Vectors-style NMF), so it is a property of the split protocol, not
one architecture, though we do not claim a discriminatively-trained network such
as 6MapNet collapses identically. Second, models graded against their own
clusters reach nDCG@10 of about 0.99, but on an external benchmark anchored to
realized like-for-like replacement transfers, no learned representation we test
beats a raw-statistics nearest-neighbour on all-candidates retrieval, the
primary metric -- a null the benchmark is powered to reject a 0.007-MAP
advantage against. This holds for a faithful Player Vectors reproduction, a
purpose-trained metric learner, and genuinely orthogonal richer inputs
(event-transformer, VAEP, and 360 spatial features), and replicates across two
disjoint player pools; within-role retrieval is near-floor and underpowered, so
we confine claims to all-candidates. We also correct a dependence error in the
field's significance testing: about 1363 queries collapse to roughly 22 position
components (effective count about 8), so per-query tests understate variance and
the naive p is about ten times too small.
\end{abstract}

\section{Introduction}
\label{sec:intro}

Player-similarity search (``find a midfielder like X''), style fingerprinting,
and representation learning over event and spatial data are all active areas
of football analytics. A less examined question is how these representations
are \emph{evaluated}. We find two pervasive pitfalls, and in both cases the
problem is that information about the relevance signal leaks into the very
model being graded.

The first pitfall is circular relevance. Similarity quality is judged against
labels derived from the model itself, such as k-means archetypes or role
clusters, or else against anecdotes. A model graded on its own partition can
score nearly perfectly; our own legacy text benchmark, whose relevance labels
are the model's k-means clusters, yields nDCG@10 of about 0.99, while the same
model fails any external test. The second pitfall is match-context leakage.
Identity and consistency are scored by splitting a player's events at random
within a tournament, so events from the same match land in both the query and
the gallery. The model can then recognize ``same game'' rather than enduring
style.

We instantiate an external, leakage-free protocol for each pitfall and report
what learned representations actually achieve once the leak is removed. The
two leaks share a single underlying cause, which we develop in
Section~\ref{sec:unifying}: the evaluation's assumed independence between the
relevance signal and the model under test is violated. They break that
independence in different places, so each needs its own fix.

We treat the two contributions asymmetrically, and say so plainly. The label-leakage
study (Pitfall~1) is the paper's primary contribution: it ships a full
external benchmark, many baselines including a faithful published method, and
a cross-pool generalization result. The feature-leakage study (Pitfall~2) is a
focused cautionary result on one task and one dataset family, with a
correspondingly narrower claim. The two-by-two view in
Section~\ref{sec:unifying} organizes the framework; it is not a claim that the
two results carry equal evidential weight.

\paragraph{Our contributions are:}
\begin{enumerate}[leftmargin=*]
\item An evaluation-leakage framework that distinguishes \emph{label leakage}
(the relevance labels are produced by the model under test) from
\emph{feature leakage} (the query and gallery share nuisance input features),
and shows that both are special cases of a violated independence assumption.
\item ScoutBench, an external, transfer-anchored similarity-retrieval
benchmark, together with the resulting negative result: no learned
representation we test, including a faithful reproduction of Player Vectors,
beats a raw-statistics nearest-neighbour baseline, and the finding replicates
across two disjoint player pools.
\item A cross-match player-identity protocol and the collapse it exposes: a
set-transformer drops from 0.878 to 0.188 Top-1 once the query and gallery are
drawn from disjoint matches, with a same-match control isolating match-context.
\item A dependence and significance correction. The field's default per-query
bootstrap treats clustered queries as independent and inflates significance by
roughly an order of magnitude; we develop the corrected block-bootstrap and
cluster-level tests as a first-class methodological contribution.
\end{enumerate}

\section{Related Work}
\label{sec:related}

The methodological survey by Davis et al.~\cite{davis2024methodology} catalogs
the evaluation protocols used across sports analytics, including the
consecutive-season nearest-neighbour test for player similarity, and the
broader pitfalls that attend them. It supports the framing we adopt here: in
practice, player-similarity and representation work is validated either
circularly, against labels the model itself produces, or qualitatively,
through scouting case studies.

The pattern is visible across the representative learned representations.
Player Vectors~\cite{decroos2019playervectors} builds per-type non-negative
matrix factorizations over action-location heatmaps and is evaluated through
consecutive-season player-identity retrieval, where the labels are objective,
supplemented by qualitative scouting case studies; it does not appeal to
external transfer ground truth. RisingBALLER~\cite{adjileye2024risingballer},
a players-as-tokens transformer foundation model, is evaluated quantitatively
on next-match-statistics prediction, while its similar-player claims remain
qualitative. SoccerMix~\cite{decroos2020soccermix}, a soft-clustering mixture
model for action and style, is evaluated qualitatively, as is
Football2Vec~\cite{magdaci2021football2vec}, a 32-dimensional Word2Vec-style
player embedding assessed through qualitative cosine nearest neighbours. A
recent spatial-statistics approach, the spatial similarity index of
G\'omez-Rubio et al.~\cite{gomezrubio2024spatial} (Lee's statistic and Moran's
$I$ over kernel-smoothed position heatmaps), is likewise demonstrated only
through illustrative case studies, with no external relevance ground truth and
no released benchmark.
Across this body of work, similarity is therefore validated either against the
model's own structure or by inspection, never against an external outcome.

Our work sits in the genre of method-level reality checks. A Metric Learning
Reality Check~\cite{musgrave2020realitycheck} showed that reported gains in
deep metric learning largely evaporate under fair protocols, and the
tabular-learning literature has repeatedly found that gradient-boosted trees
match or beat deep models on typical tabular
data~\cite{shwartzziv2022tabular,grinsztajn2022whytrees}. Dimensional collapse
in contrastive learning~\cite{jing2022dimensionalcollapse} offers a plausible
mechanism for why a learned low-dimensional projection can be lossy relative
to its raw inputs. Our raw-versus-learned result is the football-similarity
instance of these findings.

External-outcome evaluation does exist in football analytics, but it targets
quantities other than similarity. Market-value
prediction~\cite{marketvalue2023} forecasts a scalar value rather than a
similarity ranking, and transfer-impact
forecasting~\cite{dinsdale2022transferportal} predicts the effect of a move.
Neither uses realized like-for-like replacements as a similarity-retrieval
ground truth, which is the precise gap ScoutBench fills. The closest neighbour
to our identity task is 6MapNet~\cite{kim2021sixmapnet}, a triplet network for
player identity from tracking data. It splits its data by sampling player
phases per entity into train, validation, and test, rather than by match, so
phases from a single match can fall into both query and gallery; it may
therefore be vulnerable to the same match-context leakage we study, and we
treat it as a candidate to re-evaluate under a match-disjoint split rather than
as a safe reference.

\section{A Unifying View: Evaluation Leakage}
\label{sec:unifying}

Both pitfalls are, at bottom, one failure: the evaluation's assumed
independence between the relevance signal and the model under test is
violated. In label leakage, the relevance labels are produced by the model
itself, through its own clusters or archetypes, so the model is graded against
its own output and scores near-perfectly but meaninglessly. In feature
leakage, the query and the gallery share input features drawn from the same
match, so the model can match shared nuisance context instead of the target
property of enduring identity or style. Table~\ref{tab:twobytwo} arranges the
two mechanisms against the two task families.

\begin{table}[t]
\centering
\caption{The two leakage mechanisms by task family, with the fix that restores
the broken independence assumption. The two results are not of equal weight:
the label-leakage row is the paper's primary contribution, the feature-leakage
cell a focused cautionary result.}
\label{tab:twobytwo}
\begin{tabular}{lll}
\toprule
Leakage source & Similarity retrieval & Identity / consistency \\
\midrule
Label leakage & self-derived archetype relevance, & n/a (identity labels \\
              & nDCG about 0.99; fix: external & are objective) \\
              & transfer anchor (ScoutBench) & \\
\addlinespace
Feature leakage & mitigated (candidates are & within-tournament split-half, \\
                & other players) & 0.88; fix: match-disjoint split \\
\bottomrule
\end{tabular}
\end{table}

These are two distinct contamination mechanisms, not one mechanism described
twice. Label leakage violates label-model independence, because the relevance
labels are generated by the model under test. Feature leakage violates
query-gallery independence, because the query and gallery share nuisance
match-context beyond the target property, even though the identity label
itself is objective and independent of the model. Valid retrieval evaluation
requires both conditions to hold; each pitfall breaks a different one, and
each fix restores the broken one, an external label anchor for the first and a
match-disjoint split for the second. The shared lesson is not a single
mechanism but a discipline: identify every path from the relevance signal to
the representation, and sever it.

\section{Pitfall 1: Circular Relevance, and ScoutBench}
\label{sec:scoutbench}

\subsection{The benchmark}
\label{sec:scoutbench-task}

ScoutBench Task B grounds similarity in realized like-for-like replacement
transfers drawn from CC0 transfermarkt data. When a club loses player $X$ at
sub-position $P$ and signs player $Y$ at $P$ in the same or next window, the
pair $(X, Y)$ is taken as a \emph{silver} similarity-positive, that is, a label
that is objective and outcome-derived but noisy rather than a clean
gold-standard judgment. Each query player is ranked against all others by
embedding cosine similarity. We report MAP, Hit@$k$, and MRR over all
candidates as the primary metric, and over same-sub-position candidates as a
stratified secondary metric. There are 1363 query players drawn from 5868
pairs, against a gallery of 5668 players derived from StatsBomb open data and
used for research only.

\subsection{Significance done right}
\label{sec:scoutbench-stats}

The query-to-relevant graph collapses into 22 connected components, roughly
one per sub-position, so the queries are not independent. We therefore use a
block bootstrap that resamples whole components ($B = 10000$) and, separately,
a conservative cluster-level paired $t$-test over the 22 component means, with
a Holm correction across the comparison family. The field's default per-query
bootstrap, which we argue is wrong here, treats the 1363 dependent queries as
independent. It understates variance by roughly a factor of 1.3 and makes the
naive $p$-value about ten times too small, that is, it inflates significance by
roughly an order of magnitude. We treat this correction as a methodological
contribution in its own right. We also note that the query-weighted block
bootstrap and the component-equal-weighted cluster $t$-test estimate different
targets, the effective component count being about 8, so we report both rather
than presenting one as a conservative version of the other.

\subsection{Results}
\label{sec:scoutbench-results}

Table~\ref{tab:leaderboard} reports the leaderboard. All-candidates MAP
(ALL-MAP) is the primary metric; same-position MAP (SP-MAP) is the stratified
within-role test, which we show below (Sec.~\ref{sec:scoutbench-scope}) is
near-floor and underpowered. The headline result, on all-candidates, is that no
learned representation beats the raw-statistics nearest-neighbour baseline,
while every method clears random.

\begin{table}[t]
\centering
\caption{ScoutBench Task B leaderboard. SP-MAP is same-sub-position MAP;
ALL-MAP is all-candidates MAP (the primary metric). Higher is better.}
\label{tab:leaderboard}
\begin{tabular}{lcc}
\toprule
Method & SP-MAP & ALL-MAP \\
\midrule
fbref percentile & 0.0514 & 0.0082 \\
pca(card) & 0.0501 & 0.0195 \\
raw\_card cosine & 0.0498 & 0.0194 \\
TF-IDF text & 0.0469 & 0.0165 \\
Player-Vectors (faithful, Decroos \& Davis heatmap-NMF) & 0.0464 & 0.0172 \\
NMF (over card features) & 0.0450 & 0.0135 \\
v11 (contrastive embed) & 0.0447 & 0.0162 \\
card+VAEP & 0.0440 & 0.0088 \\
random & 0.0369 & 0.0030 \\
\bottomrule
\end{tabular}
\end{table}

A genuine signal does exist on all-candidates retrieval. The raw card beats
random by 0.0163 ALL-MAP, which is significant under both the block bootstrap
($p = 0.0002$) and the conservative component $t$-test ($p = 0.0001$). On
same-position retrieval the signal is weaker: it is significant under the block
bootstrap ($p = 0.0002$) but not under the component $t$-test ($p = 0.15$),
which is consistent with the position prior dominating within a role, since a
random vector already reaches SP-Hit@10 of about 0.18. The robust
signal-exists claim is therefore the all-candidates one, and the same-position
regime sits near the floor.

The null is well-powered on the primary metric and near-floor on the stratified
one, which we establish by measuring the minimum detectable effect (MDE) of the
conservative component $t$-test directly from the per-component variance (power
0.8, $\alpha = 0.05$; \texttt{scoutbench\_power.py}). On all-candidates the MDE
is about 0.007 MAP, below the best-minus-random spread of 0.0163, so the test
resolves spread-sized effects: it detects raw beating random (component
difference $+0.0097$, $p \approx 0$) and would register a
learned-representation advantage of 0.007 MAP or more, yet the learned methods
sit at or below raw (raw minus v11 $+0.0008$, $p = 0.49$; raw minus Player
Vectors $-0.0027$, $p = 0.47$). The all-candidates null is therefore genuine,
not a failure of power. On same-position the MDE is about 0.017 MAP, exceeding
the entire 0.0145 spread, so the within-role test cannot resolve even raw
beating random ($p = 0.15$) because the position prior dominates. We accordingly
draw no conclusions from same-position differences, including the apparent
raw-minus-v11 gap; they fall below the label's resolution. This is itself the
lesson: circular self-labels manufacture an effect (nDCG of about 0.99) far
above this floor, which is precisely why they look near-perfect.

The off-objective contrastive embedding does not beat raw. Here ``off-objective''
means the embedding was trained on a different objective than replacement
retrieval. On all-candidates, the primary and well-powered metric, v11 sits at
or below raw (raw minus v11 $+0.0008$, component $t$-test $p = 0.49$), a powered
null. The nominal raw-minus-v11 same-position gap of 0.0051 is significant only
under the query-weighted block bootstrap ($p$ about 0.02); it does not survive
the Holm correction and is not robust to the component $t$-test ($p = 0.36$),
and given the same-position MDE of about 0.017 it falls below the benchmark's
resolution, so we rest no claim on it.

The faithful, published Player-Vectors baseline also fails to beat raw, which
is the key external-baseline check. Our reproduction of Player
Vectors~\cite{decroos2019playervectors} builds 50-by-50 action-location
heatmaps per type (pass start and end, and shot, cross, and dribble starts),
applies per-type NMF, draws purely on on-ball events, covers the full gallery
of 5668 players, and uses none of our card features. It scores SP-MAP 0.0464
and ALL-MAP 0.0172. It is not significantly different from the raw card: the
raw-minus-PV gap is 0.0034 SP-MAP (block $p = 0.15$, component $t$
$p = 0.90$) and 0.0022 ALL-MAP ($p = 0.15$). The canonical published learned
representation thus ties or slightly trails a raw-statistics
nearest-neighbour, which moves the claim from ``our models lose'' to ``learned
representations, including the published Player-Vectors method, do not beat
raw.'' A secondary NMF over the card features behaves the same way, scoring
SP-MAP 0.0450, below raw on the query-weighted estimand ($p = 0.003$) but not
robust to the component-level multiplicity correction.

Purpose-trained, player-disjoint metric learning does not significantly beat
raw, which is the decisive, non-tautological test. We train a projection
directly on the replacement objective using a multi-positive InfoNCE loss that
masks each player's other true replacements, since the replacement graph is
many-to-many and a naive one-positive-per-row loss would create contradictory
labels. Evaluated on player-disjoint held-out queries across three seeds, the
same-position MAP gap of the metric learner minus raw is +0.0008, +0.0006, and
$-0.0027$, all with $p > 0.37$, while the all-candidates gap is +0.0016,
+0.0024, and +0.0009, consistently small and positive but never significant
($p$ between 0.19 and 0.68). A predeclared TOST equivalence test, that is, two
one-sided tests at a $\pm 0.005$ MAP margin, concludes statistical equivalence
to raw in one of the three seeds and is inconclusive in the other two, where
the 95\% confidence intervals marginally exceed the margin, for example
$[-0.006, +0.007]$. The honest statement is that there is no improvement over
raw in any seed, and equivalence is only partially established: a metric
learned directly on these features does not beat raw cosine, though we cannot
fully rule out a sub-0.005-MAP gain.

Richer, orthogonal inputs also do not beat raw on the primary metric, which is
the strongest test of the claim. Every representation above only reshapes the
88-dimensional card, so we additionally test inputs the card does not contain: a
256-dimensional event-transformer embedding over SPADL event sequences (full
coverage, with similarity-structure Pearson correlation of about 0.01 against
the card, hence genuinely orthogonal), a 3-dimensional VAEP representation, and
a StatsBomb-360 spatial fingerprint on a 923-player tournament subset. Fused
with the card via training-free cosine nearest-neighbour, none beats the raw
card on all-candidates (card+event SP-MAP 0.0503 against raw's 0.0498; card+VAEP
and card+event+VAEP at or below raw; card+360, on its subset, 0.0936 against
raw's 0.1007), and a metric MLP trained over the fused inputs (player-disjoint,
three seeds) likewise does not (SP-MAP gap $+0.004 \pm 0.004$, all seeds
non-significant). The one positive lead, reported on the same footing as every
other claim, is the event embedding alone within-position: SP-MAP \textbf{0.0703
against raw's 0.0498} (block $p = 0.0002$), the largest within-role effect we
observe. We do not count it as a win, but for exactly the reason the within-role
regime decides nothing: it fails the component $t$-test ($p = 0.11$) and does
not transfer to all-candidates (0.0093 against 0.0194), held to the identical
near-floor, underpowered standard (same-position MDE about 0.017) we apply to
raw-versus-random and raw-versus-v11, not a stricter one. Event sequences are
thus the most promising direction for within-role discrimination and the obvious
next experiment, not noise to be dismissed; they simply do not clear the bar on
the metric the benchmark can adjudicate.

Finally, the circular alternative makes the contrast concrete. The same v11
model scores nDCG@10 of about 0.99 on a benchmark whose relevance is its own
k-means labels. External grading reverses the verdict.

\subsection{Honest scope of the negative result}
\label{sec:scoutbench-scope}

The absolute numbers are low. Every method retrieves the realized replacement
rarely, with same-position Hit@10 of about 0.27, and a random vector already
reaches 0.175 there because the position prior dominates. The contribution is
the ordering and the external protocol, not deployable accuracy. The safe claim
is that no learned representation we test beats a raw-statistics
nearest-neighbour on all-candidates external transfer-anchored similarity --- a
null the benchmark is powered to overturn (MDE about 0.007, below the method
spread), not a large-margin positive result; we confine it to all-candidates
because same-position is near-floor and underpowered. The faithful
Player-Vectors reproduction makes this a claim about a published learned method,
not only about our own models. Two baselines remain for future work:
RisingBALLER, a transformer foundation model requiring days of training, and
6MapNet's exact triplet convolutional network under a cross-match split, which
needs tracking data we lack --- we test the static heatmap family it builds on
under P2 below.

\subsection{Generalization across player pools}
\label{sec:scoutbench-generalization}

We re-run Task B independently on two disjoint pools, split by the league from
which the card was built. Pool A is the top-five men's club leagues, with 3049
players and 756 queries; Pool B is international tournaments together with
women's competitions and the ISL and others, with 2619 players and 479
queries. The conclusion replicates in both. In Pool A no learned
representation beats raw: raw scores 0.0639 against v11's 0.0637
(non-significant) and Player-Vectors' 0.0598 (non-significant). In Pool B the
same holds: raw scores 0.0841 against v11's 0.0648 and Player-Vectors' 0.0677,
with neither learned method ahead. Raw beats random in both pools, by 0.0182 in
Pool A (block $p < 0.001$ and component $t$ $p = 0.003$, the cleanest
signal-exists result) and by 0.0222 in Pool B (block $p = 0.007$). The finding
is therefore not an artifact of a specific pool. Absolute MAPs are higher
within each homogeneous pool, since replacement retrieval is easier among more
similar candidates.

\section{Pitfall 2: Within-Tournament Match-Context Leakage in 360 Identity}
\label{sec:p2}

\subsection{Task and protocol}
\label{sec:p2-task}

The task is player-identity retrieval from anonymized 360 freeze-frames. For
each on-ball event we form up to 23 player point-tokens, each carrying an
$(x, y)$ location and an actor, teammate, or opponent role. A set-transformer
is trained contrastively and evaluated on held-out players, with 701 players
for training and 350 for test, over 528{,}059 events drawn from World Cup 2022,
Euro 2024, and Euro 2020 across 166 matches. The model is identical across all
conditions; only the query and target split changes.

\begin{table}[t]
\centering
\caption{Within-tournament player-identity Top-1 and MRR by split condition,
with 95\% player-level bootstrap confidence intervals. The split-half and
cross-match intervals are disjoint.}
\label{tab:p2}
\begin{tabular}{lcc}
\toprule
Condition & Top-1 [95\% CI] & MRR \\
\midrule
split-half (random, standard) & 0.878 [0.846, 0.905] & 0.924 \\
same-match control (disjoint halves of one match) & 0.889 [0.859, 0.917] & 0.932 \\
cross-match (disjoint matches, leakage-free) & 0.188 [0.153, 0.224] & 0.326 \\
cross-match, no teammate tokens & 0.137 [0.107, 0.171] & 0.274 \\
reference: cross-tournament / random & 0.178 / 0.007 & --- \\
\bottomrule
\end{tabular}
\end{table}

\subsection{Results}
\label{sec:p2-results}

Table~\ref{tab:p2} reports the conditions. The confidence intervals are
player-level bootstraps that resample held-out players. The split-half and
cross-match intervals are disjoint, so the collapse is unambiguous. It is not a
single-pool artifact: the macro-average that weights each of the three
tournaments equally matches the pooled number (split-half 0.88, cross-match
0.19), and every tournament collapses, with cross-match per-tournament Top-1 of
0.24, 0.15, and 0.19 against split-half values of 0.95, 0.86, and 0.83.

The collapse itself is large. Top-1 falls from 0.878 under split-half to 0.188
under cross-match, a drop of 0.69, at comparable event volume, about 137
against 141 events per side, so it is not a data artifact. Within-match
information is near-sufficient: the same-match positive control scores 0.889,
so when the query and target share a match, identity is nearly perfect, and
when they do not, it collapses. One caveat applies, however. The control halves
a single match, leaving about 71 events per side, so it is not volume-matched
to the split-half and cross-match conditions at about 140 events per side; it
isolates the presence of match-shared information rather than a
volume-controlled effect. The defensible conclusion is that random
within-profile splitting substantially exploits match-specific information, and
we cannot cleanly separate match identity from team, lineup, and
role-within-that-match from genuine short-term individual behaviour.

A small leakage-free signal nonetheless survives. Cross-match Top-1 of 0.188
greatly exceeds the random baseline of 0.007, dropping to 0.137 once teammate
tokens are removed. Some individual or structural signal persists, but its
interpretation as individual style is bounded by the confounds named below.

The collapse is not specific to our architecture, with a bounded generalization
claim. \emph{Static} heatmap representations collapse the same way under a
match-disjoint split. A parameter-free multi-channel heatmap (actor location,
teammate density, and opponent density) drops from 0.806 (95\% CI [0.768,
0.838]) under random split-half to 0.113 ([0.083, 0.146]) cross-match; a Player
Vectors-style NMF reduction drops from 0.595 ([0.549, 0.639]) to 0.123 ([0.093,
0.155]). Both show disjoint split-half and cross-match intervals and a high
same-match control. The leak is therefore a property of the within-tournament
random-split protocol, not of any single architecture: it strikes a
parameter-free histogram and an NMF reduction exactly as it strikes our trained
set-transformer. An important scope caveat is that these are static, untrained
representations. We do \emph{not} claim a discriminatively-trained triplet
convolutional network such as 6MapNet collapses identically, since a learned
model could in principle acquire match-invariant features a static histogram
cannot. What we show is that the protocol leaks for the heatmap input family on
which published methods rest --- Player Vectors directly, and 6MapNet's heatmap
inputs --- while 6MapNet's exact network needs continuous tracking data we lack
and remains untested.

The implication should be stated carefully, and we do not accuse the
literature. Within-tournament random split-half is a tempting,
seemingly-natural protocol that would inflate identity by roughly 4.7 to 6
times through match-context leakage. The established literature, however, does
not obviously use random within-tournament split-half. Player
Vectors~\cite{decroos2019playervectors} evaluates consecutive-season
consistency, which is match-disjoint with respect to exact same-match overlap,
though team and role context can still recur across seasons.
6MapNet~\cite{kim2021sixmapnet}, the closest neighbour, samples player phases,
potentially several from one match, into train, validation, and test rather
than splitting by match; and the static heatmap input family it builds on
collapses under our match-disjoint split (above), so a phase-sampling protocol
like 6MapNet's is exposed in principle. Whether its trained network actually
collapses is untested --- it needs tracking data we lack --- so we flag it as a
candidate to re-evaluate, not a proven case and not a safe reference. Our result is
therefore a cautionary protocol finding together with a cheaper safe
alternative: cross-match-within-season splitting recovers match-disjointness
without requiring multiple seasons, and a same-match positive control should
accompany any identity claim.

Two residual threats deserve naming, though neither is fatal. First, train-time
leakage: the encoder was trained with random-half, same-match positives, so it
may have learned to exploit match context, and a model trained with cross-match
positives might recover signal. We therefore claim only that, as trained and
evaluated, the 0.88 figure reflects match context, not that individual style is
unlearnable from 360 data. Second, the no-teammate ablation at 0.137 also
removes about half the input tokens, confounding model capacity with leakage
type; a cleaner test would permute teammate identities rather than delete
tokens, which we leave to future work.

\section{Leakage-Free Protocols}
\label{sec:protocols}

We package the two fixes as reusable protocols. ScoutBench Task B is an
external, transfer-anchored similarity-retrieval evaluation with an open
evaluator, a block bootstrap paired with a cluster $t$-test and Holm
correction, and all-candidates retrieval as the primary metric. The cross-match
identity protocol is a match-disjoint query and target split, volume-comparable
to split-half at about 140 events per side, accompanied by a same-match
positive control. Both ship with code and corrected baselines, and both are
cheap to adopt.

\section{Recommendations}
\label{sec:recommendations}

Three recommendations follow directly. First, never grade similarity against a
model's own labels; use external or outcome anchors, and always report the
random-baseline floor. Second, for identity and consistency, split by match or
by season, never randomly within a competition, and include a same-match
control. Third, treat queries as clustered into position components, use
cluster-level inference, and report absolute floors, effect sizes, and the
test's minimum detectable effect rather than $p$-values alone --- a null is only
informative if the benchmark could have detected a real advantage of the
relevant size (here, well-powered on all-candidates, near-floor within-role).

\section{Threats to Validity and Limitations}
\label{sec:threats}

The replacement labels are silver: realized replacements encode budget,
availability, and positional need, not pure similarity. We checked their
validity directly. Realized replacements sit at a mean similarity-percentile of
0.549 among same-position players, with a sub-position-clustered confidence
interval of $[0.521, 0.585]$ that excludes the 0.50 null, and replacement
cosine of 0.449 against a random-same-position 0.404. The label therefore
carries genuine but weak stylistic similarity signal beyond position; it is
noisy rather than pure noise. This coheres the benchmark: a weak label signal
explains why all methods beat random yet cluster near the floor and none
separates cleanly, since the raw card captures the linearly-available
similarity and the residual is mostly non-stylistic transfer noise.

The evaluation draws on a single dataset family, StatsBomb open and Wyscout
public, which is non-commercial. Cross-pool generalization is shown, since the
result replicates on top-five men's leagues and on tournaments, women's, and
other competitions (Section~\ref{sec:scoutbench-generalization}), but
cross-dataset generalization on a fully independent provider remains future
work. The learned-model coverage on ScoutBench includes a faithful published
Player-Vectors baseline that ties raw, v11, a purpose-trained metric learner,
NMF, and card+VAEP, all of which fail to beat raw; RisingBALLER as a
transformer foundation model remains future work. On the 360 task, training
with cross-match positives might learn a less-leaky encoder and partially
recover signal, which is an open question rather than a refutation.

\section{Conclusion}
\label{sec:conclusion}

Two widely-used styles of player-representation evaluation grade the model
against information that the model can see. When relevance labels come from the
model's own clusters, scores saturate near 0.99 and mean nothing; when query
and gallery come from the same match, an identity score of 0.878 collapses to
0.188 once the matches are made disjoint. In both cases the failure is the
same: a path runs from the relevance signal back into the representation under
test. The fixes differ only in where that path is cut, an external transfer
anchor for label leakage and a match-disjoint split for feature leakage, and
both are cheap to adopt. The discipline that ties them together is a single
instruction we offer as the paper's takeaway: identify every path from the
relevance signal to the representation, and sever it. Under that discipline, no
learned representation we test beats a raw-statistics nearest-neighbour on
all-candidates external transfer-anchored similarity --- the published Player
Vectors method included, and on a benchmark powered to have detected such an
advantage --- which relocates the open problem from better embeddings to better,
externally-grounded labels.

\section*{Reproducibility}
\label{sec:repro}

The ScoutBench code lives in \texttt{football\_embed/evaluation/}, comprising
\texttt{scoutbench.py}, \texttt{scoutbench\_significance.py},
\texttt{scoutbench\_blockboot.py} (block bootstrap, component $t$-test, and
Holm correction), \texttt{scoutbench\_extended.py},
\texttt{scoutbench\_metric\_baseline.py} (multi-positive loss and TOST),
\texttt{scoutbench\_label\_validity.py},
\texttt{scoutbench\_generalization.py} (cross-pool), and
\texttt{scoutbench\_power.py} (minimum detectable effect / sensitivity), with the
Player-Vectors builder at
\texttt{football\_embed/data/player\_vectors\_nmf.py}; results are
written to \texttt{data/processed/benchmark/*.json}. The 360 protocol lives in
\texttt{sb360\_crossmatch\_killtest.py} (split-half, cross-match, same-match
control, and no-teammate conditions) and
\texttt{sb360\_6mapnet\_crossmatch.py} (the static heatmap family), with data in
\texttt{sb360\_sets\_matchkeyed.npz} (offline-rebuilt and validated) and
results in \texttt{sb360\_crossmatch\_killtest.json}. The card matrix is
sanitized (clipped to $\pm 50$ and passed through nan-to-num), and aggregate
metrics are confirmed unchanged; similarity matrices are verified finite and
bounded ($|\cos| \le 1.0$ by direct checks). A spurious BLAS matmul runtime
warning appears on some macOS Arrow and Accelerate builds but does not affect
results, and is suppressed at the call site where this was verified.

\bibliographystyle{plain}
\bibliography{references}

\end{document}
